\begin{table}[htbp]
\centering
\small
\caption{生成器嵌套选参：全数据内层留出潜空间评分（nat/款）}
\label{tab:q3_select}
\resizebox{\linewidth}{!}{%
\begin{tabular}{lrrrrrr}
\hline
协方差 & k & κ & 全数据内层 L & 小样本风格内层 L & 外层选中次数 & 全数据选定 \\
\hline
famod & 15 & 5.00 & -242.33 & -253.87 & 0 & 是 \\
famod & 10 & 2.00 & -242.70 & -251.44 & 24 & 否 \\
famod & 10 & 5.00 & -242.98 & -251.80 & 1 & 否 \\
famod & 20 & 5.00 & -243.72 & -258.26 & 0 & 否 \\
famod & 25 & 5.00 & -246.69 & -263.09 & 0 & 否 \\
fa & 15 & 5.00 & -246.73 & -258.33 & 0 & 否 \\
fa & 20 & 5.00 & -247.30 & -262.09 & 0 & 否 \\
fa & 10 & 5.00 & -248.16 & -257.06 & 0 & 否 \\
fa & 25 & 5.00 & -249.31 & -265.31 & 0 & 否 \\
famod & 5 & 5.00 & -250.42 & -255.08 & 0 & 否 \\
famod & 30 & 5.00 & -251.23 & -267.44 & 0 & 否 \\
fa & 30 & 5.00 & -252.80 & -268.86 & 0 & 否 \\
famod & 2 & 5.00 & -259.05 & -262.41 & 0 & 否 \\
fa & 5 & 5.00 & -259.57 & -265.14 & 0 & 否 \\
lw & — & 5.00 & -267.87 & -286.96 & 0 & 否 \\
fa & 2 & 5.00 & -271.31 & -276.30 & 0 & 否 \\
diag & — & 5.00 & -298.49 & -304.69 & 0 & 否 \\
\hline
\end{tabular}}
\par\footnotesize 注：全数据选定 famod、k=15、κ=5；列出 κ 取选定值的设定与外层被选中的设定；完整 96 个设定见 tab\_q3\_select.csv。小样本风格 L 只报告，不参与选参。
\end{table}
